\documentclass[11pt]{article}
\usepackage[margin=1in]{geometry}
\usepackage{amsmath,amssymb,hyperref}
\newcommand{\Tr}{\operatorname{Tr}}
\title{Collective focusing in two dimensions:\\the central problem and verified reductions}
\author{Shuxing Yang}
\date{11 October 2026}
\begin{document}
\maketitle
\paragraph{Status.}
The complete two-dimensional sharp Lieb--Thirring problem is unresolved
by this work. The reductions below use existing spectral theorems and
standard NLS symmetries. The tensor-product bound is a direct application
of a classical sharp inequality; no historical novelty is asserted.
The accompanying numerical searches found no counterexample.

\section{The precise target}
For the two-dimensional common-density system
\[
 i\partial_tu_j+\Delta u_j+\rho u_j=0,\qquad
 \rho=\sum_j|u_j|^2,\qquad
 E=\tfrac12\sum_j\|\nabla u_j\|_2^2-\tfrac14\int\rho^2,
\]
the eigenvalues of the full Gram matrix are conserved occupations.
Let $-\Delta Q+Q-Q^3=0$ be the scalar ground-state equation,
$M_c=\|Q\|_2^2$, and $C_1=2/M_c$.
The question is whether, for every positive finite-trace operator
$\Gamma$ with finite kinetic energy,
\begin{equation}\label{target}
 \|\rho_\Gamma\|_2^2\le C_1\|\Gamma\|_{\rm op}\Tr(-\Delta\Gamma).
\end{equation}
It is equivalent by spectral duality to
\[
 \Tr(-\Delta-V)_-\le \frac{1}{2M_c}\int_{\mathbb R^2}V_+^2.
\]
If true, \eqref{target} yields
\[
 E\ge\tfrac12(1-\|\Gamma\|_{\rm op}/M_c)\Tr(-\Delta\Gamma).
\]
For a fixed normalized occupation spectrum $p$, the total-mass threshold
would therefore be $M_*(p)=M_c/p_{\max}$. Below this value all
finite-component energy-space solutions are global and uniformly bounded;
above it negative-energy finite-variance data exist. The latter follows
by concentrating the largest mode and putting the others on disjoint
broad supports. No assertion is made that all superthreshold data collapse.

\section{Every failure can be sought among equal occupations}
Set $m=\|\Gamma\|_{\rm op}$, $P_s={\bf1}_{(s,\infty)}(\Gamma)$ and
$D(P)=C_1\Tr(-\Delta P)-\|\rho_P\|_2^2$. The exact identity
\begin{align*}
 C_1m\Tr(-\Delta\Gamma)-\|\rho_\Gamma\|_2^2
 &=m\int_0^mD(P_s)\,ds\\
 &\quad+\tfrac12\int_0^m\int_0^m
       \|\rho_{P_s}-\rho_{P_t}\|_2^2\,ds\,dt
\end{align*}
follows from $\Gamma=\int_0^mP_s\,ds$ and the Hilbert-space variance
identity. The classical non-sharp kinetic inequality ensures integrability.
Each $P_s$ has finite rank. Thus a strict counterexample always yields
a finite-rank projection counterexample, with all occupied modes weighted
equally. This reduction does not bound its rank.

\section{A failure would produce an exact collective standing wave and collapse}
Write $L_N$ for the best constant in the first-$N$ spectral sum.
Frank--Gontier--Lewin's finite-rank optimizer, bubble-decomposition and
Euler--Lagrange theorems apply at $d=2,\gamma=1$.
If $L_N>L_1=1/(2M_c)$ first occurs at $N=N_*$, minimality excludes
splitting of maximizing sequences into smaller-rank optimizers.
There is a normalized optimizer, real analytic and exponentially decaying,
with
\[
 \int V_*^2=1,\qquad
 V_*=\frac1{2L_{N_*}}\sum_{j=1}^{N_*}|\phi_j|^2,\qquad
 (-\Delta-V_*)\phi_j=-e_j\phi_j,\qquad
 \langle\phi_i,\phi_j\rangle=\delta_{ij}.
\]
There may be additional unoccupied negative eigenvalues; they must not
be excluded by assumption.

Let $\beta=1/(2L_{N_*})<M_c$ and $\psi_j=\sqrt{\beta}\phi_j$. Then
\begin{equation}\label{standing}
 -\Delta\psi_j-\left(\sum_k|\psi_k|^2\right)\psi_j=-e_j\psi_j,\qquad
 \langle\psi_i,\psi_j\rangle=\beta\delta_{ij}.
\end{equation}
Summing the eigenvalue equations gives
$\sum_j\|\nabla\psi_j\|_2^2=1/2$, $\int\rho_\Psi^2=1$, hence $E(\Psi)=0$.
Conversely any nonzero localized solution of \eqref{standing} with
equal occupation $\beta<M_c$ disproves \eqref{target}:
the Pohozaev identity gives $E(\Psi)=0$, and normalization of the
orbitals yields the quotient $2/\beta>2/M_c$.

Choose real eigenfunctions, let $h=T-t>0$, and set
\[
 u_j(t,x)=h^{-1}
 \exp\!\left(-\frac{i|x|^2}{4h}+\frac{ie_j}{h}\right)\psi_j(x/h).
\]
Pseudoconformal invariance, or direct differentiation, gives an exact
NLS solution with Gram matrix $\beta I$ and
\[
 \sum_j\|\nabla u_j(t)\|_2^2
 =h^{-2}\sum_j\|\nabla\psi_j\|_2^2+
   \tfrac14\sum_j\||x|\psi_j\|_2^2.
\]
Its energy is $\frac18\sum_j\||x|\psi_j\|_2^2$.
It collapses although every occupation is below $M_c$.
These implications identify the unresolved existence question;
they do not construct such a subcritical standing wave.

\section{A global exclusion of complete product occupations}
The classical sharp one-dimensional $\gamma=3/2$ inequality is
\[
 \Tr(-\partial_x^2-V)_-^{3/2}\le\tfrac3{16}\int V_+^2.
\]
For a positive finite-rank operator $A$, write
$B_A=\int\rho_A^2$ and $T_A=\Tr(-\partial_x^2A)$.
Variational duality and Schatten H\"older yield, for $a>0$,
\[
 aB_A-T_A\le
 (\Tr A^3)^{1/3}(3/16)^{2/3}a^{4/3}B_A^{2/3}.
\]
Optimization in $a$ gives $B_A^2\le(\Tr A^3)T_A/3$.
For contractions $0\le A,B\le1$, therefore,
\[
 \int\rho_{A\otimes B}^2
 \le\tfrac13\sqrt{(\Tr A)(\Tr B)T_AT_B}
 \le\tfrac16\Tr[-\Delta(A\otimes B)].
\]
Convex combinations obey the same estimate by Jensen's inequality.
In particular, every complete rectangular occupation
$\{f_i(x)g_j(y):1\le i\le m,\ 1\le j\le n\}$
has kinetic quotient at most $1/6$, with no shape or rank restriction.
The constant is attained at rank one by
$\frac12\operatorname{sech}x\,\operatorname{sech}y$ and approached
at larger rectangular ranks by separated copies.

The strict inequality $1/6<C_1$ has an elementary witness.
For $q(x)=\operatorname{sech}|x|/\sqrt{2\pi\log2}$ in two dimensions,
\[
 \frac{\int q^4}{\int|\nabla q|^2}
 =\frac{4\ell-1}{2\pi\ell(2\ell+1)}>\frac16,\qquad \ell=\log2.
\]
Indeed $69/100<\ell<7/10$ and $\pi<22/7$ give the strict bound
$1.76/10.56=1/6$. The logarithm estimates follow from
$\log2=2\sum_{k\ge0}[(2k+1)3^{2k+1}]^{-1}$.
Nonrectangular occupation sets and general nonseparable projections
are not covered.

\section{Numerical exploration and the remaining mathematical task}
The recorded periodic Gaussian-well potential is
\[
 V(x)=A\sum_{z\in\mathbb Z^2}
 e^{-|x-2\pi z|^2/(2\sigma^2)}+c,
 \quad A=2.9414883,\quad \sigma=0.73521703,\quad c/A=0.01513357.
\]
Its computed spectral ratio is approximately $0.0425618561$,
below the scalar target $1/(2M_c)\simeq0.0427317684$ by $0.398\%$.
Plane-wave cutoffs and Bloch meshes were increased; the recorded digits
were stable. Fourier and spatial evaluations of the potential cost
also agree. These are floating-point checks, not rigorous enclosures.
No continuum upper bound or global optimality follows.

The open task is to construct or exclude the equal-occupation collective
structures in \eqref{standing}, without radiality, fixed-rank or
single-core assumptions. The existing reductions and product bound
do not accomplish this. Full derivations, source checks, code and
numerical results are included in the accompanying working files.

\begin{thebibliography}{9}
\bibitem{FGL}
R.~L. Frank, D.~Gontier and M.~Lewin,
\emph{Optimizers for the finite-rank Lieb--Thirring inequality},
arXiv:2109.05984v2 (2023), Theorem 1, Corollary 2 and Theorem 3.
\url{https://arxiv.org/abs/2109.05984}
\bibitem{FGL2}
R.~L. Frank, D.~Gontier and M.~Lewin,
\emph{The nonlinear Schr\"odinger equation for orthonormal functions II.
Application to Lieb--Thirring inequalities}, CMP (2021).
\url{https://arxiv.org/abs/2002.04964}
\bibitem{LW}
A.~Laptev and T.~Weidl,
\emph{Sharp Lieb--Thirring inequalities in high dimensions},
Acta Math.\ (2000).
\url{https://arxiv.org/abs/math-ph/9903007}
\end{thebibliography}
\end{document}

